
This chapter assesses the legal, social, ethical and professional dimensions of building and releasing an automatic annotator. The detailed ethics record is in Appendix A.

\section{Legal considerations}

\textbf{Data licensing and attribution.} The source dataset is published under CC-BY 4.0 (Wang et al., 2025), which permits reuse and derivative works with attribution. Attribution is given throughout this dissertation, on the validation website, and in the released repository; the automatic annotations are a derivative work distributed under the same attribution terms.

\textbf{Model licences as a design constraint.} Licensing was an engineering requirement, and in one case it shaped the pipeline: Depth Anything v2 is Apache-2.0 only in its Small variant, the larger ones being non-commercial, so Small was preferred, a choice ablation A8 later justified independently on accuracy (Appendix D.5). SAM2 and Grounding DINO are Apache-2.0. The benchmark detector's YOLO training uses the \texttt{ultralytics} library, which is AGPL-3.0: compatible with academic research use, but a commercial deployment of that component would need a licence review. The note is recorded so a future user does not inherit the question unknowingly.

\textbf{Data protection.} Some frames contain identifiable people, making them personal data under UK law (Data Protection Act 2018). The project processes them for research under the research provisions of that Act; the CC-BY licence governs reuse of the material and does not by itself settle the data-protection question, which is why the basis is stated separately. It applies data minimisation to what it republishes: faces are anonymised in all published figures and on the validation website, and items whose judgement anonymisation would compromise are removed rather than shown. No data is collected from human participants anywhere in the project. No unauthorised access occurs anywhere in the project, so the Computer Misuse Act 1990 is noted only for completeness.

\section{Social considerations}

\textbf{Automation of annotation work.} The project automates paid human work. The framing the measurements support is this: at this scale the manual process produced sparse, inconsistent labels at a cost of nine annotators, and the realistic effect of automation is not displacing a profession but changing the human role, from labelling every pair to reviewing a measurable flagged minority (\S{}4.7). The same shift makes dataset construction affordable for groups that could not fund manual annotation at all, since the pipeline runs on a consumer 6 GB GPU.

\textbf{Downstream robot behaviour.} Annotation quality propagates: a robot planner consuming wrong spatial relations can act wrongly in physical space. This is precisely why the project's evaluation centres on audited precision, abstention instead of guessing, and per-failure attribution, and why Chapter 6 distinguishes labels that are \textit{correct} from labels that are \textit{human-like}. Trust in automated labels should be calibrated by exactly the kind of evidence this dissertation supplies, not assumed.

\textbf{Bias.} The geometric rules themselves carry no demographic component: they compute from positions and extents. The perception stack, however, is learned, and detection quality for the \texttt{human} class cannot be assumed uniform across people; published detectors have documented performance disparities. In this dataset the people are members of the collecting research group, so the question is not testable here, and it is recorded as a deployment consideration, not a resolved issue.

\section{Ethical considerations}

The ethical surface has two parts, each with a concrete safeguard. First, secondary use of images containing identifiable people: face anonymisation in everything republished, and item removal where anonymisation would obscure the object under judgement rather than showing it. \textbf{No data is collected from human participants at any point}, so the work is secondary analysis throughout and the Secondary Data Checklist is the applicable route (Appendix A). Second, research integrity: predictions were registered before the benchmark run and reported as they fell, one confirmed, one refuted and one left unresolved once a better-controlled replication shrank the margin it rested on (\S{}6.6); a withdrawn hypothesis remains in the text (\S{}6.4), and two built refinements are reported as measured and declined (Appendix D.4).

\section{Professional considerations}

The work follows the BCS Code of Conduct (BCS, 2022): public interest, in the privacy safeguards above; professional competence and integrity, in claims bounded by measurements and limitations stated in \S{}7.6; and duty to the profession, in methods and failures documented so others can build on both. In the engineering itself: version control with a clean history, a unit-and-invariant suite run before every change ships, every threshold in one configuration file, seeded reproducible runs (Appendix B), and licence-compliant use of third-party models and data. The reproducibility package is a deliverable with the same status as the results, because a result that cannot be re-run is an anecdote.

\section{Summary}

Automating annotation raises obligations in four registers, and this chapter has set out how each is discharged. Legally, the dataset's CC-BY licence permits the derivative labels, model licences shaped one design decision, and the frames containing identifiable people are personal data handled under the research provisions of the Data Protection Act 2018. Socially, the realistic effect is a changed human role rather than a displaced profession, with the residual review cost measured rather than asserted. Ethically, the safeguards are face anonymisation, secondary analysis only with no human participants, and an audit trail that records superseded results. Professionally, the work follows the BCS Code of Conduct in its claims and in its engineering practice.

Chapter 9 closes the dissertation against its objectives and research questions, states the contributions, and sets out what is left undone.